\clearpage\fancyhead[L]{\small ARITHMETISCH GEOMETRISCHES MITTEL}
\pdfbookmark[0]{ARITHMETISCH GEOMETRISCHES MITTEL}{work-native370}
\SourceLeaf{370}
\selectlanguage{english}
\begin{center}
\textbf{NACHLASS.}\par
\textbf{[ARITHMETISCH GEOMETRISCHES MITTEL.]}\par
\textbf{PARS I.}\par
\textbf{DE ORIGINE PROPRIETATIBUSQUE}\\\textbf{GENERALIBUS NUMERORUM}\par
\textbf{MEDIORUM ARITHM. GEOMETRICORUM.}\par
\rule{0.09\linewidth}{0.4pt}\par
\end{center}
\AGMArticle{1}
Sint $\left\{\begin{array}{l}a,\allowbreak \ a',\allowbreak \ a'',\allowbreak \ a'''\ .\,\allowbreak .\,\allowbreak .\,\allowbreak .\\ b,\allowbreak \ b',\allowbreak \ b'',\allowbreak \ b'''\ .\,\allowbreak .\,\allowbreak .\,\allowbreak .\end{array}\right\}$ duae progressiones quantitatum ea lege formatae, ut
quilibet ipsarum termini correspondentes sint \textit{media} inter terminos antecedentes,
et quidem termini progressionis superioris media arithmetica, progressionis inferioris geometrica, puta
\[
\begin{aligned}&a'=\tfrac12(a+b),\quad b'=\surd ab,\\[6pt]&a''=\tfrac12(a'+b'),\quad b''=\surd a'b',\\[6pt]&a'''=\tfrac12(a''+b''),\quad b'''=\surd a''b''\ \mathrm{etc.}\end{aligned}
\]
Supponemus autem, ipsos $a$, $b$ esse reales positivos, et pro radicibus quadraticis
ubique accipi valores positivos; quo pacto progressiones quousque libuerit produci
poterunt, omnes ipsarum termini erunt plene determinati valoresque positivos reales nanciscentur. Porro prima hic se fronte offerunt observationes sequentes:
\par
I. Si $a=b$, omnes utriusque seriei termini erunt $=a=b$.
\par
II. Si vero $a$, $b$ sunt inaequales, erit $(a'-b')(a'+b')=\tfrac14(a-b)^2$, unde
concluditur $b'<a'$, et perinde erit $b''<a''$, $b'''<a'''$ etc., i.\,e. quivis terminus seriei inferioris minor erit quam correspondens superioris. Quocirca in hoc casu
supponemus, esse etiam $b<a$.
\par
III. Eadem suppositione erit $a'<a$, $b'>b$; $a''<a'$, $b''>b'$ etc.; progressio
itaque superior continuo decrescit, inferior continuo crescit; hinc manifestum est,
utramque habere limitem; hi limites commode exprimuntur per $a^{\infty}$, $b^{\infty}$.
\par
IV. Denique ex $\frac{a'-b'}{a-b}=\frac{(a-b)}{4(a'+b')}=\frac{a-b}{2(a+b)+4b'}$ sequitur $a'-b'<\tfrac12(a-b)$,
eodemque modo erit $a''-b''<\tfrac12(a'-b')$ etc. Hinc concluditur, $a-b$, $a'-b'$,
\SourceLeaf{371}
$a''-b''$, $a'''-b'''$ etc. constituere progressionem continuo decrescentem atque ipsius
limitem esse $=0$. Hinc $a^{\infty}=b^{\infty}$, i.\,e. progressio superior et inferior eundem
limitem habebunt, quo illa semper manet maior, haec minor.
\par
Hunc limitem vocamus \textit{numerum medium arithmetico-geometricum inter $a$ et $b$},
et per $\mathrm{M}(a,b)$ designamus.
\AGMArticle{2}
Radices aequationis $xx-2ax+bb=0$ erunt reales positivi, siquidem
$a\geqq b$; medium arithmeticum inter has radices erit $a$, geometricum $b$; designata
itaque una radice (et quidem maiore si sunt inaequales) per ${}^{\prime}\!a$, altera per ${}^{\prime}\!b$, poterit ${}^{\prime}\!a$ spectari tamquam terminus progressionis superioris terminum $a$ praecedens, eodemque modo ${}^{\prime}\!b$ tamquam terminus progressionis inferioris ante $b$. Similiter designando
\begingroup\setlength{\arraycolsep}{2pt}
\[
\begin{array}{@{}lllll@{}}
\mathrm{aequationis}&xx-2\,{}^{\prime}\!ax+{}^{\prime}\!b\,{}^{\prime}\!b=0&\mathrm{radicem\ maiorem\ per}&{}^{\prime\prime}\!a&\mathrm{minorem\ per}\ {}^{\prime\prime}\!b\\[6pt]
&xx-2\,{}^{\prime\prime}\!ax+{}^{\prime\prime}\!b\,{}^{\prime\prime}\!b=0&&{}^{\prime\prime\prime}\!a&\hfill{}^{\prime\prime\prime}\!b\\[6pt]
&xx-2\,{}^{\prime\prime\prime}\!ax+{}^{\prime\prime\prime}\!b\,{}^{\prime\prime\prime}\!b=0&&{}^{\prime\prime\prime\prime}\!a&\hfill{}^{\prime\prime\prime\prime}\!b
\end{array}
\]\endgroup
poterunt ${}^{\prime}\!a$, ${}^{\prime\prime}\!a$, ${}^{\prime\prime\prime}\!a$ etc. spectari tamquam continuatio progressionis superioris versus laevam, atque ${}^{\prime}\!b$, ${}^{\prime\prime}\!b$, ${}^{\prime\prime\prime}\!b$ etc. tamquam continuatio progressionis inferioris, ita
ut iam habeantur duae progressiones utrimque in infinitum continuabiles
\[
\begin{array}{@{}ll@{}}
.\,.\,.\,.\ {}^{\prime\prime\prime\prime}\!a,\;{}^{\prime\prime\prime}\!a,\;{}^{\prime\prime}\!a,\;{}^{\prime}\!a,\; a,\; a',\; a'',\; a''',\; a''''\ .\,.\,.\,.\,.&\mathrm{(I)}\\[6pt]
.\,.\,.\,.\ {}^{\prime\prime\prime\prime}\!b,\;{}^{\prime\prime\prime}\!b,\;{}^{\prime\prime}\!b,\;{}^{\prime}\!b,\; b,\; b',\; b'',\; b''',\; b''''\ .\,.\,.\,.\,.&\mathrm{(II)}
\end{array}
\]
\par Quivis itaque terminus progressionis (I) erit maior quam correspondens seriei (II); series illa a laeva ad dextram continuo decrescit, a dextra ad laevam
crescit: haec a laeva ad dextram continuo crescit sensuque contrario decrescit.
Versus dextram utraque series eundem limitem habet; versus laevam autem (I)
super omnes limites crescit, (II) habet limitem $0$ (nisi omnes utriusque progressionis termini sunt aequales). Nam ${}^{\prime}\!a=a+\surd(aa-bb)$; ${}^{\prime}\!b=a-\surd(aa-bb)$;
hinc ${}^{\prime}\!a\,{}^{\prime}\!a-{}^{\prime}\!b\,{}^{\prime}\!b=4a\surd(aa-bb)>4(aa-bb)$, similiterque ${}^{\prime\prime}\!a\,\allowbreak {}^{\prime\prime}\!a-{}^{\prime\prime}\!b\,\allowbreak {}^{\prime\prime}\!b>4({}^{\prime}\!a\,\allowbreak {}^{\prime}\!a-{}^{\prime}\!b\,\allowbreak {}^{\prime}\!b)$
etc., unde patet, seriem $aa-bb$, ${}^{\prime}\!a\,{}^{\prime}\!a-{}^{\prime}\!b\,{}^{\prime}\!b$, ${}^{\prime\prime}\!a\,{}^{\prime\prime}\!a-{}^{\prime\prime}\!b\,{}^{\prime\prime}\!b$ etc. et proin etiam hanc
$a$, ${}^{\prime}\!a$, ${}^{\prime\prime}\!a$ etc. quemvis limitem superare posse; $\frac{{}^{\prime}\!b}{b}=\frac{b}{{}^{\prime}\!a}<\frac{b}{a}$ similiterque $\frac{{}^{\prime\prime}\!b}{{}^{\prime}\!b}<\frac{{}^{\prime}\!b}{{}^{\prime}\!a}$ etc.
i.\,e. $\frac{{}^{n+1}\!b}{{}^{n}\!b}$ infra quemvis limitem deprimi potest augendo ipsum $n$, adeoque limes
seriei $b$, ${}^{\prime}\!b$, ${}^{\prime\prime}\!b$ etc. $=0$.
